% =====================================================================
% Abstract — closed-loop sentence(s). Scope-locked to sections (OC-2/3, no
% claim stronger than §7.4/§7.6). Does NOT use "comprehensive"; detour worded
% as "low" with distribution; P2 = interface-robust.
% =====================================================================
\begin{abstract}
We study how the \emph{output interface} of language-grounded routing should represent traffic
constraints, and introduce \textbf{STAG-Route}, which grounds natural-language constraints into typed
sparse anchors rather than dense per-edge weights. On a held-out, 150-case, method-agnostic SUMO stress
test in which the unconstrained shortest path always violates the constraint, STAG-Route attains high
conditional constraint satisfaction ($0.96$) at low excess detour (median $0.0$s; 96\% of cases $\le 60$s),
ranking first among nine methods on both satisfaction and travel time. In contrast, a dense edge-weight
grounder collapses into a fixed-capacity, scenario-invariant enumeration that fails as $|E|$ grows---failing
both constraint satisfaction ($0.02$) and efficiency---and a near-oracle numeric (sensor-based) baseline fails \emph{structurally} on
the multi-edge buckets that define the task, a failure invariant to sensor coverage. Ablations show the
typed-anchor interface is robust across model families (Qwen2.5, Llama-3.2) and held-out topologies at the
road level, while a verifier provides schema/reachability guarantees over an already-clean,
hallucination-free model. We further characterize the method's frontier: grounding is invariant to scale,
topology, and language, but degrades under \emph{compositional} out-of-distribution stress (multiple
simultaneous conflicting events), which we report and reconcile with the closed-loop result.
\end{abstract}
